\scanpage{8}
\hfill\handwritten{\unsure{Ag|AG}}

% slanted
\margin{\handwritten{Matrices de Coxeter\\ \unsure{Matrice|Matrices} de Cartan\\ Systèmes de racines\\ groupes engendrés par des \unsure{réflexions}\\ diagrammes de Dynkin}}

\begin{center}GRIE ET ALGOUPES DE LEBRES.\hfill N\textsuperscript{o}383\end{center}

\begin{center}CHAPITRE III.\end{center}

\begin{center}SYSTINES DE RACEMES (Etat 3).\end{center}

\begin{center}\mbox{-{}-{}-{}-{}-{}-{}-{}-{}-{}-{}-{}-}\end{center}

\begin{center}\uline{Sommaire}\end{center}

\begin{center}\S~1. \uuline{Résultats préliminaires}.\end{center}

1. Lemmes sur les chambres \dotfill{} 1

2. Lemmes sur les espaces hilbertiens réels de dimension finie \dotfill{} 3

3. Lemmes de Géométrie plane \dotfill{} 5

\begin{center}\S~2. \uuline{Groupes de Coxeter}.\end{center}

1. Définition des groupes de Coxeter \dotfill{} 1

2. Représentation directe d’un groupe de Coxeter \dotfill{} 3

3. Représentation contragrédiente d’un groupe de Coxeter \dotfill{} 5

4. Démonstration du théorème 2 \dotfill{} 8

5. La réunion des chambres \dotfill{} 15

6. Critère de finitude d’un groupe de Coxeter \dotfill{} 16

\begin{center}\S~3. \uuline{Groupes engendrés par des réflexions}.\end{center}

1. Groupes finis engendrés par des réflexions vectorielles \dotfill{} 1

2. Groupes affines engendrés par des réflexions \dotfill{} 6

3. Points spéciaux \dotfill{} 13

\begin{center}\S~4. \uuline{Systèmes de racines}.\end{center}

1. Définition d’un système de racines \dotfill{} 1

2. Relations entre deux racines \dotfill{} 5

3. Systèmes de racines réduits \dotfill{} 9

4. Chambres. Racines positives. Racines simples \dotfill{} 11

5. Plus grande racine \dotfill{} 13

6. Ensembles clos de racines \dotfill{} 13

7. Groupes commutatifs libres définis par un système de racines \dotfill{} 17
% archive sheet number: 7
